\chapter{Modular Arithmetic: The Number System Cryptography Lives In}
\label{ch:modular}

\section{Clock arithmetic, taken seriously}

You already compute modulo twelve every day: four hours after ten o'clock is
two o'clock. Wrap-around arithmetic --- add, overflow the dial, keep the
remainder --- is the entire idea of \emph{modular arithmetic}. Cryptography's
only twist is the size of the clock: Ed25519's dial has
\[
p = 2^{255} - 19
\]
hours on it --- a 77-digit prime. Everything else --- the wrap-around, the
remainder-taking, the algebra --- is the twelve-hour clock you already know.

Formally, we write $a \equiv b \pmod{n}$ when $n$ divides $a - b$, and we
collect all integers with the same remainder into one object: the world
$\Zmod{n}$ has exactly $n$ elements, $\{0, 1, \dots, n-1\}$, with addition
and multiplication that wrap. In Lean/Mathlib this world is a first-class
type, and computation in it is exactly what you would hope:

\begin{lstlisting}[language=Lean]
#eval (7 + 8 : ZMod 12)    -- 3     (fifteen o'clock is three o'clock)
#eval (5 * 9 : ZMod 12)    -- 9
#eval (3 - 7 : ZMod 12)    -- 8     (subtraction wraps -- no truncation!)

def p : Nat := 2^255 - 19
#eval (2^255 : ZMod p)     -- 19    (of course: 2^255 = p + 19)
\end{lstlisting}

Note the third line with relief: unlike \lean{Nat}, subtraction in
\lean{ZMod n} is a total, well-behaved inverse of addition. Every element has
a negative. In algebra terms, $\Zmod{n}$ is a \emph{commutative ring} --- and
Lean knows it, so the \lean{ring} tactic from Chapter~\ref{ch:automation}
works there natively.

\section{Division: where primes enter}

Rings give us $+$, $-$, $\times$. Cryptography also needs $\div$ --- and
division is where the modulus stops being a size parameter and starts being a
design decision. Dividing by $a$ means multiplying by some $a^{-1}$ with
$a \cdot a^{-1} = 1$. Does such an inverse exist? On the twelve-hour clock,
try to invert $4$: the multiples of $4$ are $4, 8, 0, 4, 8, 0, \dots$ ---
the value $1$ never appears. $4$ has no inverse mod $12$, because $4$ and
$12$ share the factor~$4$.

\begin{bigidea}
In $\Zmod{n}$, the element $a$ is invertible exactly when
$\gcd(a, n) = 1$. So if $n = p$ is \textbf{prime}, \emph{every} nonzero
element is invertible: you can divide freely, and $\Zmod{p}$ earns the title
of \textbf{field}, written $\Fp$. Fields are the arithmetic paradise where
linear algebra, polynomial factoring, and elliptic-curve geometry all work.
This --- and only this --- is why cryptographic moduli are prime: primality
is the entry ticket to division.
\end{bigidea}

\begin{center}
\begin{tikzpicture}[scale=0.62,every node/.style={font=\small}]
  % Z/12 clock: multiples of 4 stuck in a cycle
  \begin{scope}
    \draw[ink2,thick] (0,0) circle (2.2);
    \foreach \i in {0,...,11} \node[color=ink] at ({90-\i*30}:1.8) {\i};
    \foreach \s/\t in {0/4, 4/8, 8/0}
      \draw[-{Stealth},accent,thick] ({90-\s*30}:2.55) arc ({90-\s*30}:{90-\t*30+8}:2.55);
    \node[color=accent,align=center] at (0,-3.4) {$\Zmod{12}$: stepping by $4$\\ visits only $\{0,4,8\}$ --- never $1$};
  \end{scope}
  % Z/11 clock: multiples of 4 reach 1
  \begin{scope}[xshift=9.5cm]
    \draw[ink2,thick] (0,0) circle (2.2);
    \foreach \i in {0,...,10} \node[color=ink] at ({90-\i*32.72}:1.8) {\i};
    \foreach \s/\t in {0/4, 4/8, 8/1}
      \draw[-{Stealth},proven,thick] ({90-\s*32.72}:2.55) arc ({90-\s*32.72}:{90-\t*32.72+8}:2.55);
    \node[color=proven,align=center] at (0,-3.4) {$\Zmod{11}$: stepping by $4$\\ reaches $1$ in three steps: $4^{-1}=3$};
  \end{scope}
\end{tikzpicture}
\end{center}

How do you \emph{find} $a^{-1}$ mod $p$? Two classical answers, both of which
appear in real Ed25519 code: the extended Euclidean algorithm, and Fermat's
little theorem, which says $a^{p-1} \equiv 1 \pmod p$ for $a \not\equiv 0$
--- hence $a^{p-2}$ \emph{is} the inverse. The dalek library computes
inverses as $a^{p-2}$ with a hand-crafted chain of $254$ squarings and $11$
multiplications; one of the proofs you will meet in Chapter~\ref{ch:field}
verifies precisely that this chain computes what Fermat promises.

\begin{worked}{inverting $19$ modulo the 77-digit prime, entirely by hand}
A 77-digit modulus does not put pen-and-paper out of business --- watch.
We compute $19^{-1} \bmod p$ for the real $p = 2^{255}-19$, exactly, in
five lines of Euclid. The only preparation is one residue:
\emph{what is $p$ mod $19$?} By Fermat's little theorem in the
\emph{small} field $\Zmod{19}$: $2^{18} \equiv 1$, and
$255 = 14 \cdot 18 + 3$, so
\[
2^{255} \equiv (2^{18})^{14} \cdot 2^3 \equiv 8 \pmod{19}
\qquad\Longrightarrow\qquad
p = 2^{255} - 19 \equiv 8 - 0 \equiv 8 \pmod{19}.
\]
So $p = 19k + 8$ for the (77-digit, but never-written-out) integer
$k = (p-8)/19$. Now run Euclid on $(p, 19)$, keeping remainders only:
\[
p = 19k + 8, \qquad 19 = 2\cdot 8 + 3, \qquad 8 = 2\cdot 3 + 2,
\qquad 3 = 1\cdot 2 + 1 .
\]
The gcd is $1$ (of course --- $p$ is prime). Back-substitute, bottom to
top, expressing $1$ in terms of each earlier pair:
\[
\begin{array}{lcl}
1 &=& 3 - 2\\
  &=& 3 - (8 - 2\cdot 3) \;=\; 3\cdot 3 - 8\\
  &=& 3\,(19 - 2\cdot 8) - 8 \;=\; 3\cdot 19 - 7\cdot 8\\
  &=& 3\cdot 19 - 7\,(p - 19k) \;=\; (3 + 7k)\cdot 19 \;-\; 7p .
\end{array}
\]
Read off the answer: $19 \cdot (3 + 7k) = 1 + 7p \equiv 1 \pmod p$, so
\[
19^{-1} \bmod p \;=\; 3 + 7k
\;=\; 3 + 7\cdot\frac{p - 8}{19} .
\]
This is an \emph{exact, closed-form} answer about the real prime, computed
with numbers that never exceeded two digits --- the size of $p$ entered
only through one residue. Two lessons. First, modular arithmetic is
constantly this kind: gigantic numbers handled through small
representatives --- the same move the limb representation of
Chapter~\ref{ch:denotation} industrializes. Second, note what the
computation \emph{cost}: five divisions. Fermat's route costs $254$
squarings of 77-digit numbers. Why does the real code use Fermat anyway?
Because Euclid's division sequence \emph{depends on the input value} ---
its running time leaks secrets through timing --- while the Fermat chain
executes identically for every input. Constant-time discipline pays a
couple hundred multiplications of overhead for silence; you will meet
this trade at every layer of the pyramid.
\end{worked}

\section{\texorpdfstring{Why $2^{255}-19$?}{Why 2**255-19?} A prime chosen for machines}

Any large prime makes a field. Why this one? Because arithmetic mod $p$ is
computed by \emph{machines with 64-bit words}, and $p = 2^{255}-19$ is shaped
for them:

\begin{itemize}[leftmargin=1.4em]
\item \textbf{Reduction is a multiply-by-19.} Numbers at or beyond $2^{255}$
  overflow the dial; but since $2^{255} \equiv 19 \pmod p$, any overflow bit
  re-enters the sum carrying a factor of just $19$. Reducing mod $p$ costs
  one small multiplication --- no long division, ever.
\item \textbf{255 bits split evenly into five limbs of 51.} A 64-bit word
  holding a 51-bit limb leaves 13 bits of headroom for carries to accumulate
  lazily --- the delayed-carry style from Chapter~\ref{ch:why}, now by design
  rather than accident. The bound $a_i < 2^{51+\varepsilon}$ will follow us
  through Chapters~\ref{ch:denotation} and~\ref{ch:field}.
\item \textbf{It sits just below a power of two}, so 255-bit values fit in 32
  bytes with one bit to spare --- and Ed25519 spends that spare bit storing a
  point's sign. Engineering all the way down.
\end{itemize}

\begin{worked}{the multiply-by-19 reduction, performed on real overflow}
Do the reduction the code does, by hand, on the real modulus. The
identity: $2^{255} = p + 19 \equiv 19 \pmod p$. Now take a genuinely
overflowing value --- say $2^{260}$, five bits past the edge, the kind of
thing a lazy carry chain produces. Split off the overflow and fold:
\[
2^{260} \;=\; 2^{5} \cdot 2^{255}
\;\equiv\; 2^{5} \cdot 19
\;=\; 608 \pmod{p}.
\]
One multiplication by a two-digit constant replaced a division by a
77-digit prime. Once more, with a mixed value: reduce
$x = 2^{256} + 2^{255} + 7$:
\[
x = (2 + 1)\cdot 2^{255} + 7 \;\equiv\; 3 \cdot 19 + 7 \;=\; 64 \pmod p .
\]
And the general recipe, exactly as the code implements it: write
$x = x_{\mathrm{lo}} + 2^{255} x_{\mathrm{hi}}$ (a bit-split, free in
hardware); return $x_{\mathrm{lo}} + 19\, x_{\mathrm{hi}}$. Estimate the
shrinkage like an engineer: if $x < 2^{300}$, then
$x_{\mathrm{hi}} < 2^{45}$ and the result is
$< 2^{255} + 19\cdot 2^{45} < 2^{256}$ --- one pass tames $45$ excess bits,
and a second pass lands below $2^{256}$-ish for good. This
split-multiply-add \emph{is} the entire reduction machinery of every
curve25519 implementation on earth; you have now executed it. When
Chapter~\ref{ch:field} verifies \code{reduce}, the theorem is the
congruence line above, quantified over every $x$ the bounds admit.
\end{worked}

The Pasta curves verified in the companion projects (Pallas/Vesta, used in
zero-knowledge proof systems) choose their $\approx 2^{254}$ primes by a
different criterion --- friendliness to fast Fourier transforms --- and
represent elements in \emph{Montgomery form}, a representation trick we will
touch in Chapter~\ref{ch:denotation}. Different shapes, same theory: pick a
prime whose field your machine can love.

\begin{worked}{shadow arithmetic --- auditing big computations on small clocks}
Your grandmother's bookkeeper knew a modular-arithmetic trick worth
resurrecting: \emph{casting out nines}. Because $10 \equiv 1 \pmod 9$,
a number is congruent mod $9$ to its digit sum --- so any claimed sum or
product of large numbers can be spot-checked by comparing digit-sum
shadows. The general principle: \textbf{any exact integer identity
survives reduction mod anything}, so a cheap clock can audit an
expensive computation. Watch it catch the classic transcription error
--- two leading digits transposed:
\[
123456789 \times 987654321 \;\overset{?}{=}\;
\underline{21}1932631112635269 .
\]
Shadow mod $9$: digit-sum$(123456789) = 45 \equiv 0$, so the product's
shadow must be $0 \cdot (\text{anything}) = 0$; the claimed result's
digit sum is $63 \equiv 0$ --- \emph{passes}. It would: transpositions
never change a digit sum, which is exactly why bookkeepers paired nines
with a second clock. Mod $11$ ($10 \equiv -1$, so take the
\emph{alternating} digit sum from the right): each factor
$\equiv 5 \pmod{11}$, so the product must be
$5 \cdot 5 = 25 \equiv 3 \pmod{11}$; the claimed result's alternating
sum gives $1 \not\equiv 3$ --- \textbf{caught}. (The true product is
$121932631112635269$; the transposition slipped past one clock and not
the other, which is the design lesson: independent shadows multiply
their power, $1$-in-$9$ and $1$-in-$11$ escape odds compounding to
$1$-in-$99$.)

Why this card lives in this book: the verified field code is audited by
the \emph{same architecture}. The denotation bridge of
Chapter~\ref{ch:denotation} checks limb computations by their shadow in
$\Zmod{p}$ --- one gigantic well-chosen clock instead of two small ones
--- and the kernel plays the incorruptible bookkeeper. And when
\emph{you} debug a failing proof about a $77$-digit computation,
shadowing both sides mod $9$ or mod $97$ by hand remains the fastest
way to find out \emph{which} side is lying. (Exercise 6.6 hands you a
broken carry chain to convict this way.)
\end{worked}

\section{Modular arithmetic in Lean, hands on}

Statements about $\Fp$ are ordinary Lean theorems, and the automation you
already own applies:

\begin{lstlisting}[language=Lean]
-- ring identities hold verbatim in ZMod n:
example (x y : ZMod 12) : (x + y)^2 = x^2 + 2*x*y + y^2 := by ring

-- small finite worlds: check every case
example : ∀ x : ZMod 12, 4 * x ≠ 1 := by decide

-- inverses exist in prime fields (Mathlib knows ZMod 11 is a field):
example (a : ZMod 11) (h : a ≠ 0) : a * a⁻¹ = 1 :=
  ZMod.mul_inv_cancel_of_ne_zero h
\end{lstlisting}

What about the crown fact, $2^{255} \equiv 19$? For \emph{concrete numeric}
statements at 77-digit scale, tactic choice suddenly matters: \lean{decide}
would ask the kernel to grind case analysis it cannot afford, while
\lean{norm_num} and reflexivity-by-computation on efficient numerals handle
it comfortably. Verifying real cryptography is partly a \emph{performance
engineering} discipline --- Chapter~\ref{ch:prime} turns this observation
into a principle when the question becomes primality itself.

\begin{pitfall}
In \lean{ZMod n}, the numeral \lean{57896044618658...} and the numeral
\lean{19} can be \emph{the same element}. Equality of elements is not
equality of the numerals you typed. When a surprising \lean{rfl} succeeds ---
or two ``different'' constants turn out equal --- remember you are on the
clock face, not the number line. This is a feature: the entire verification
story of Chapter~\ref{ch:denotation} consists of choosing, deliberately, when
to view a pile of bytes as an integer and when as a clock position.
\end{pitfall}

\begin{tryit}
Open \code{exercises/Ch06.lean}. You will: compute your first inverses with
\lean{\#eval}; expand $(x+y)^3$ in \lean{ZMod 7} with one tactic; show $4$
has no inverse in \lean{ZMod 12} by \lean{decide}; and prove the reduction
identity $2^{255} = 19$ in \lean{ZMod p}.
\end{tryit}

\section*{Exercises}

\exercise{Compute by hand (yes, hand): $3^{-1}$ in $\Zmod{7}$, and check that
stepping by $3$ on a 7-hour clock visits every position. Then confirm both
with \lean{\#eval}.}

\exercise{Prove in Lean: \lean{∀ x : ZMod 12, 4 * x ≠ 1}, then change $12$
to $13$ and watch the same tactic refuse --- find the witness it is
implicitly pointing at.}

\exercise{Fermat's little theorem is \lean{ZMod.pow_card_sub_one_eq_one} in
Mathlib. Use it to prove that in \lean{ZMod 11}, \lean{a\textasciicircum{}9 * a = 1} whenever
\lean{a ≠ 0}.}

\exercise{(Paper) The Ed25519 code reduces a 256-bit value $x$ by writing
$x = x_{\mathrm{low}} + 2^{255}\, x_{\mathrm{high}}$ and returning
$x_{\mathrm{low}} + 19\, x_{\mathrm{high}}$. Prove informally that this
preserves the value mod $p$, and bound how much smaller the result is. You
have just re-derived the reduction step you will verify formally in
Chapter~\ref{ch:field}.}

\exercise{(Paper) Run the worked Euclid computation for a different small
constant: find $121665^{-1} \bmod p$? No --- that one needs more lines
than a margin holds. Do the honest scaled version: compute
$p \bmod 3$ (via $2^{255} \bmod 3$) and derive $3^{-1} \bmod p$ in closed
form, as the worked example did for $19$. Check your first Euclid line by
verifying the claimed residue.}

\exercise{(Shadow audit) A limb computation claims
$(2^{51} - 19) + (2^{51} - 1) = 2^{52} - 21$. Convict or acquit using
two shadows: mod $9$ and mod $5$. (Powers of $2$ cycle mod $9$ with
period $6$ and mod $5$ with period $4$ --- Card~3 of the toolkit
generalizes.)}

\section*{Solutions and pathways}
\solutionsintro

\solhead{6.1}
\pathway ``Inverse of $3$ mod $7$'' asks: three times what is one more
than a multiple of $7$? With seven candidates, enumerate --- and while
enumerating, watch the stepping-around-the-clock picture from the chapter
figure happen.
\answer $3 \cdot 5 = 15 = 2\cdot 7 + 1$, so $3^{-1} = 5$ in $\Zmod{7}$.
Stepping by $3$ from $0$: $0 \to 3 \to 6 \to 2 \to 5 \to 1 \to 4 \to 0$
--- all seven positions visited before returning, because
$\gcd(3,7)=1$. \lean{\#eval (3⁻¹ : ZMod 7)} confirms \lean{5}. (If you
try the same walk stepping by $4$ on the 12-clock, you revisit $0$ after
three steps --- the figure's left panel --- and that early return
\emph{is} the nonexistence of the inverse.)

\solhead{6.2}
\pathway Twelve cases is \lean{decide} territory; the interest is in the
refusal after you change the modulus. A decision procedure failing on a
false statement is \emph{pointing} at something --- make it tell you what.
\answer \lean{example : ∀ x : ZMod 12, 4 * x ≠ 1 := by decide} succeeds.
With modulus $13$: \lean{decide} fails, because the statement is now
false --- $\gcd(4,13) = 1$, so an inverse exists. Find it as the worked
examples teach: $4 \cdot 10 = 40 = 3 \cdot 13 + 1$, so the witness is
$x = 10$, confirmed by \lean{\#eval (4⁻¹ : ZMod 13)} printing \lean{10}.
The general moral, third time now: a decision procedure's ``no'' is a
counterexample's ``here.''

\solhead{6.3}
\pathway Fermat gives $a^{10} = 1$ directly ($p - 1 = 10$ at $p = 11$);
the exercise's $a^9 \cdot a$ is that same fact wearing a disguise ---
undo the disguise with exponent algebra ($a^9 \cdot a = a^{10}$), then
apply the library lemma. Registering \lean{Fact (Nat.Prime 11)} is the
price of admission for the typeclass machinery.
\answer
\begin{lstlisting}[language=Lean]
instance : Fact (Nat.Prime 11) := ⟨by decide⟩

example (a : ZMod 11) (h : a ≠ 0) : a^9 * a = 1 := by
  have hp := ZMod.pow_card_sub_one_eq_one h   -- a^10 = 1
  calc a^9 * a = a^10 := by ring
       _       = 1    := by simpa using hp
\end{lstlisting}
The two-step shape --- \emph{algebraic massage by \lean{ring}, then the
library fact} --- is how nearly every ``mathematical'' step in the real
proofs goes. The inversion-chain verification of
Chapter~\ref{ch:field} is this exercise, iterated 265 times.

\solhead{6.4}
\pathway Congruence first (replace $2^{255}$ by $19$ and check the
difference is a multiple of $p$), size second (bound each summand).
\answer \emph{Congruence:} the returned value differs from $x$ by
$x_{\mathrm{high}} \cdot (2^{255} - 19) = x_{\mathrm{high}} \cdot p$, a
multiple of $p$ --- so they are congruent. (That is the whole proof; write
it as one displayed equation and admire how little was needed.)
\emph{Size:} for $x < 2^{256}$ we have $x_{\mathrm{high}} < 2$, i.e.\
$x_{\mathrm{high}} \in \{0, 1\}$, and $x_{\mathrm{low}} < 2^{255}$, so the
result is $< 2^{255} + 19 < 2^{256}$ --- and if $x_{\mathrm{high}} = 1$
the value dropped by $p - 19 \approx 2^{255}$: one pass halves the range.
The formal statement in \code{dalek-ed25519-verified} carries exactly
these two clauses --- congruence and bound --- the two-clause spec shape
of Chapter~\ref{ch:denotation}, met here in miniature.

\solhead{6.5}
\pathway Same three moves as the worked example: (i) small-field Fermat
for the residue, (ii) Euclid (here trivially short), (iii)
back-substitute into closed form.
\answer Residue: $2 \equiv -1 \pmod 3$, so
$2^{255} \equiv (-1)^{255} = -1 \equiv 2 \pmod 3$, hence
$p = 2^{255} - 19 \equiv 2 - 1 \equiv 1 \pmod 3$ (since
$19 \equiv 1 \bmod 3$). So $p = 3m + 1$ with $m = (p-1)/3$. Then
$3m = p - 1 \equiv -1$, i.e.\ $3 \cdot (-m) \equiv 1$, giving
\[
3^{-1} \bmod p \;=\; p - m \;=\; p - \frac{p-1}{3} \;=\; \frac{2p + 1}{3}.
\]
Sanity-check the closed form: $3 \cdot \frac{2p+1}{3} = 2p + 1 \equiv 1
\pmod p$ ✓. Notice the answer is one line shorter than the worked
example's --- because $p \bmod 3 = 1$ made Euclid collapse. The residue
does all the work; the size of $p$ never mattered.

\solhead{6.6}
\pathway Shadow each side independently; a mismatch on \emph{any} clock
convicts. Powers of $2$ cycle: mod $9$ with period $6$
($2^6 = 64 \equiv 1$), mod $5$ with period $4$ ($2^4 \equiv 1$) ---
reduce exponents by the period, exactly Card~3's move.
\answer Mod $9$: $51 \equiv 3 \pmod 6$ so $2^{51} \equiv 2^3 = 8$, and
$19 \equiv 1$, so the left side is
$\equiv (8 - 1) + (8 - 1) = 14 \equiv 5$.
Claimed right side: $52 \equiv 4 \pmod 6$ so $2^{52} \equiv 2^4 = 7$,
and $21 \equiv 3$, giving $7 - 3 = 4$. Shadow verdict: $5 \neq 4$ ---
\textbf{convicted}. Cross-examine mod $5$: $2^{51} \equiv 2^{3} = 3$
(period $4$), left $\equiv (3 - 4) + (3 - 1) = -1 + 2 = 1$; right:
$2^{52} \equiv 1$, $21 \equiv 1$, gives $0$. Again $1 \neq 0$ ---
convicted twice. The true sum is $2^{52} - 20$ (check its shadows:
mod $9$: $7 - 2 = 5$ ✓; mod $5$: $1 - 0 = 1$ ✓); the claim was off by
exactly one --- the species of error this book has been hunting since
Chapter~\ref{ch:why}, caught this time with two clocks and thirty
seconds.

\begin{checkpoint}
You should now be able to: compute in $\Zmod{n}$ and explain the notation;
state exactly when division works and why primality guarantees it; give two
independent reasons the constant $19$ appears throughout curve25519
codebases; and prove small modular facts in Lean with \lean{decide},
\lean{ring}, and a Mathlib lemma found by name.
\end{checkpoint}
